% Work draft: actual finite and compatible Heisenberg projections of cusp origins.
\subsection{Actual Heisenberg projections of two cubic cusp origins}

\begin{statement}{B.43}{Two actual cusp-path projections and their exact finite values}{Proved}
Let $K=\mathbb Q(\zeta)$, where $\zeta^2+\zeta+1=0$, and retain
the plus Fermat cubic and the $O=C_0=[1:-1:0]$ tangent.
The two morphisms to the tripod
\[
 f_1=\frac{(1+\zeta)X}{Y+Z},\qquad
 f_2=\frac{X}{X+\zeta^2Y}                                  \tag{B.43.1}
\]
give explicit compatible Heisenberg projections of the
actual class-two origins at $A_0$ and $C_1$, respectively.
Their interior base values are
\[
 z_1=f_1(O)=\zeta^2,\qquad
 z_2=f_2(O)=\frac1{1-\zeta^2}.                              \tag{B.43.2}
\]
The projected cusp origins are specified at every finite
level by the actual radical fibres below, including their
marked tangential fibres.  This is an exact compatible
description of these two projections, not a finite
approximation declared equal to an infinite three-adic number.

On the two homogeneous boundary action directions of B.41,
the resulting changes of cusp origins are, in the intrinsic
fibre-functor Hall commutator coordinate, $9\psi_1$ and
$54\psi_2$. Their deck-generator coordinates have the opposite signs.
These are differences with a fixed abelian origin.  The
raw projected origins are nonabelian section classes and
must not be identified with either linear cocycle before
subtracting a specified reference origin.
\end{statement}
\begin{proof}
\par\smallskip\noindent\textbf{The actual maps and tangent frames.}
The divisors of the first map and its complement are
\[
 \begin{split}
 \operatorname{div}(f_1)&=A_1+A_2-2A_0,\\
 \operatorname{div}(1-f_1)&=B_1+C_2-2A_0.
 \end{split}                                                \tag{B.43.3}
\]
Indeed $Y+Z$ is the flex tangent at $A_0$, whereas
$Y+Z+\zeta^2X$ is the secant through $A_0,B_1,C_2$.
For the second map,
\[
 \begin{split}
 \operatorname{div}(f_2)&=A_0+A_1+A_2-3C_1,\\
 \operatorname{div}(1-f_2)&=B_0+B_1+B_2-3C_1.
 \end{split}                                                \tag{B.43.4}
\]
Thus both maps are defined on the nine-punctured cubic.
The values in (B.43.2) follow by substitution.
In the source parameters $s_{A_0}=X/Z$ and $s_{C_1}=Z/X$,
the target parameter $r=1/f_i$ has the exact leading terms
\[
 r=\frac{\zeta}{3}s_{A_0}^{2}+O(s_{A_0}^{5}),\qquad
 r=-\frac13s_{C_1}^{3}+O(s_{C_1}^{6}).                       \tag{B.43.5}
\]
These follow by expanding $(1+s^3)^{1/3}$, whose coefficients
are integral at every prime other than three.  In particular
the tangential comparison retains both the ramification
indices and the leading scalar; an ordinary point-path
value alone would omit these data.

\par\smallskip\noindent\textbf{A compatible explicit Heisenberg cover.}
For $m=3^r$ choose compatible primitive roots with
$\xi_{3m}^3=\xi_m$ and $\xi_m^{m/3}=\zeta$, write
$\xi=\xi_m$, and use
\[
 u^m=t,\qquad v^m=1-t,\qquad
 w^m=E_m(u),\quad
 E_m(u)=\prod_{i=0}^{m-1}(1-\xi^iu)^i.                     \tag{B.43.6}
\]
All functions being rooted are units over the tripod, so
this is finite etale there.  At the tangent $0\!\to\!1$
choose $u=t^{1/m}$, and choose the roots $v,w$ with
Puiseux constant one.  The following are deck transformations:
\[
 a:(u,v,w)\mapsto(\xi u,v,w(1-u)/v),\qquad
 b:(u,v,w)\mapsto(u,\xi v,w).                              \tag{B.43.7}
\]
The identity which verifies the first formula is
\[
 \frac{E_m(\xi u)}{E_m(u)}
       =\left(\frac{1-u}{v}\right)^m.                     \tag{B.43.8}
\]
At $t=1$, the factor which could vanish at $u=1$ has
exponent zero, so $w$ has trivial local monodromy; this
identifies the positive loop at one with the deck label
$b$ on the chosen branch reached from $u=0$ to $u=1$.
The analogous label for the positive loop at zero is $a$.
This does not identify $b$ with the monodromy permutation
on every sheet.  Label the fibre by $g\cdot p_0$ using
left deck transformations.  Analytic continuation commutes
with these transformations, so the positive loops act by
right multiplication $R_a,R_b$.  This is the same reversal
of chronological monodromy used in B.40.
Both $a$ and $b$ have order $m$, and in the convention of
composition of deck transformations
\[
 h=[b,a]:(u,v,w)\mapsto(u,v,\xi w).                        \tag{B.43.9}
\]
The intrinsic fibre-functor Hall generator is therefore
\[
 c=[R_b,R_a]=R_{[b,a]^{-1}}=R_{h^{-1}}.                     \tag{B.43.9a}
\]
Since $h$ is central its left and right actions agree.
Its intrinsic degree-two Lie coordinate is consequently
$-c$. All Lie projections below use $c$; all radical
residue coordinates use the displayed deck generator $h$.
Consequently the monodromy group is the free class-two,
exponent-$m$ group on $a,b$, of order $m^3$.
The two local monodromies generate it, proving connectedness.
This identifies (B.43.6), with its displayed marking, with
the characteristic quotient
$F_2/(\gamma_3F_2,F_2^m)$.
It follows in particular that the reductions
$a\mapsto a$, $b\mapsto b$, $h\mapsto h$ give the compatible
tower as $m$ varies.  The arithmetic descent is the unique
one preserving the stated Puiseux roots at the rational
tangent; equivalently it is the characteristic arithmetic
quotient of the tripod.  No theorem restricted to roots of
unity of order prime to three is used here.
More explicitly, if $\sigma(\xi)=\xi^c$ and $r_c$ is the
inverse of $c$ in $\{1,\ldots,m-1\}$, set
\[
 Q_{r_c}(u)=\prod_{j=0}^{m-1}
          (1-\xi^ju)^{\lfloor r_cj/m\rfloor}.
\]
The arithmetic lift fixing the tangent roots $u,v$ is
\[
 \sigma(w)=w^{r_c}/Q_{r_c}(u).                             \tag{B.43.6a}
\]
The $m$th powers agree because reindexing the product
gives $E_{m,\xi^c}(u)=E_{m,\xi}(u)^{r_c}/Q_{r_c}(u)^m$.
The Puiseux constant one fixes the possible root of unity.
This also proves the action and composition identities,
and specifies the arithmetic comparison at every finite
level, not only at the primes used in the table.

Pulling (B.43.6) back along (B.43.1) gives explicit actual
finite etale covers of $U_3$.  They are connected: the
boundary exponent matrices displayed below have rank two
modulo three, so the maps onto the geometric pro-three
tripod group are surjective by its Frattini criterion.
The fibre at the $O$ tangent is the ordinary fibre of
(B.43.6) at $z_i$; regular units in the source parameter
have unique compatible roots after choosing their constant
roots.  The fibres at $A_0$ and $C_1$ are instead obtained
by substituting (B.43.5) in the Puiseux expansions at infinity.
These are actual fibres of the same covers.  Their Galois
actions therefore give, by functoriality of fibre functors,
the pushforwards of the actual marked cusp sections.
As compatible finite torsor comparisons they specify
the two projected nonabelian section classes exactly.

\par\smallskip\noindent\textbf{A finite Frobenius formula including the endpoints.}
Let $p\ne3$ be a rational prime and choose an embedding of
$K$ into $\mathbb F_p$.  Suppose
\[
 p\equiv1\pmod{3m},\qquad 1-\zeta^2\in(\mathbb F_p^*)^m.
                                                               \tag{B.43.10}
\]
Both $z_i$ and $1-z_i$ are then $m$th powers.  Choose
$u_i\in\mathbb F_p$ with $u_i^m=z_i$, and write
\[
 E_m(u_i)^{(p-1)/m}=\xi^{k_i},\qquad k_i\in\mathbb Z/m.
                                                               \tag{B.43.11}
\]
Equation (B.43.8) proves independence of the choice of
$u_i$: its quotient is an $m$th power because a root of
$1-z_i$ is also in $\mathbb F_p$.  The ordinary fibre
Frobenius is the central deck element $h^{k_i}$.

At infinity, put $r=1/t$.  Its leading product is
\[
 E_m(u)=c_\infty\,r^{-(m-1)/2}
                  (1+O(r^{1/m})),\qquad
 c_\infty=(-1)^{(m-1)/2}\zeta^2.                            \tag{B.43.12}
\]
Indeed $\sum i=m(m-1)/2$ and
$\sum i^2=m(m-1)(2m-1)/6$; the latter exponent of $\xi$
reduces to $\zeta^2$ for every power $m$ of three.
Condition (B.43.10) makes $c_\infty$ an $m$th power.
Moreover
\[
 3=-\zeta(1-\zeta^2)^2,
\]
so both leading scalars $\zeta/3$ and $-1/3$ in
(B.43.5) are $m$th powers.  All three Puiseux coordinates
$u,v,w$ at the two source cusps can therefore be chosen
with Frobenius-fixed leading coefficients.  Their higher
principal units have unique roots since $m$ is invertible
at $p$.  The marked cusp fibre has a Frobenius-fixed point.

In the same finite arithmetic cover, the base $O$ section
thus has image $h^{k_i}$ and the marked cusp section has
image one.  Hence its cocycle relative to the $O$ section is
\[
                   h^{-k_i}.                              \tag{B.43.13}
\]
This determines the sign without an implicit path convention.
Changing a connecting path conjugates the two images, and
their centrality makes (B.43.13) independent of that choice.
By (B.43.9a), the same relative cusp element has intrinsic
Hall coordinate $+k_i$.

The executed finite-field calculation gives the following
values; both $k_1$ and $k_2$ equal the displayed $k$ in these
cases. The final column is the actual relative cusp-origin
coordinate in the deck generator $h$; its intrinsic Hall
coordinate is the preceding column $k$.
\[
\begin{array}{c|r|r|r}
 m&p&k&-k\pmod m\\\hline
 3&73&1&2\\
 3&271&1&2\\
 9&271&4&5\\
 9&1621&1&8\\
 27&2269&15&12\\
 27&3079&17&10\\
 81&32563&68&13\\
 81&176419&30&51
\end{array}                                                  \tag{B.43.14}
\]
For every row, the program checks all $m$ choices of the
root $u_i$ and the exact identity (B.43.8).  At each listed
prime it also recomputes every lower level and checks
reduction of $k$.  Thus these are computations in actual
compatible finite covers; no infinite value is inferred
from their precision.

\par\smallskip\noindent\textbf{The integral homogeneous comparison and its limit.}
In the original basis $(f_1,f_2,A_0,A_1,A_2,B_0,B_1,B_2,C_1,C_2)$,
the two maps on geometric abelianization are
\[
 P_1=\begin{pmatrix}
 -1&0&-2&1&1&0&0&0&0&0\\
 -1&0&-2&0&0&0&1&0&0&1
 \end{pmatrix},\quad
 P_2=\begin{pmatrix}
 0&-1&1&1&1&0&0&0&-3&0\\
 0&-1&0&0&0&1&1&1&-3&0
 \end{pmatrix}.                                             \tag{B.43.15}
\]
The boundary columns are the divisors (B.43.3)--(B.43.4);
the first two columns follow uniquely from
$P_i(u)=P_i(v)=0$ and the integral gluing of B.24.
If $\Delta_1,\Delta_2$ are the two normalized homogeneous
derivations, use $b_1=b$ from (B.41.1) and
\[
 \begin{split}
 b_2={}&-[A_1,A_0]+[A_2,A_0]-[A_2,A_1]\\
       &-[B_1,B_0]+[B_2,B_0]-[B_2,B_1]\\
       &+[C_1-C_2,A_0+A_1+A_2-B_0-B_1-B_2]
 \end{split}
\]
in the construction (B.41.2)--(B.41.3).  This specifies
the second normalization without an unnamed basis choice.
The exact Lie calculation gives
\[
 \begin{array}{c|cc}
 &\Delta_1&\Delta_2\\\hline
 \operatorname{Lie}_3(P_1)\Delta_j&-9\operatorname{ad}_c P_1&0\\
 \operatorname{Lie}_3(P_2)\Delta_j&0&-54\operatorname{ad}_c P_2
 \end{array}.                                               \tag{B.43.16}
\]
Equivalently, apply $\Lambda^2P_i$ to the forced origin
change $3b_g(T_Pb-b)$ of B.41, at $P=A_0$ for $i=1$ and
$P=C_1$ for $i=2$.  The resulting two rows on the pair
of homogeneous cocycles are exactly
\[
                  (9,0),\qquad(0,54).                       \tag{B.43.17}
\]
Both comparisons are recorded with all integral coefficients.
They are injective on the two rational homogeneous directions,
and also on their torsion-free $\mathbb Z_3$ coefficient lines.
Equation (B.43.17) uses $c$; its deck-$h$ rows are $(-9,0)$
and $(0,-54)$. This conversion is forced by (B.43.9a),
not chosen to fit a Frobenius value.

There is a necessary affine qualification.  A projected
origin itself has nonzero abelian cocycle in general;
its central coordinate satisfies a cup-product equation,
not the linear cocycle equation in isolation.  For example,
at $p=73$ the two abelian Kummer coordinates vanish modulo
three, as do $\kappa(\zeta)$ and $\kappa(1-\zeta)$.
Its Frattini class in the free arithmetic group $P$ is
therefore zero, so every linear $H^1(P,\mathbb F_3)$
class vanishes there.  Nonetheless (B.43.14) has $k=1$.
The nonzero value is a secondary nonabelian contribution;
it cannot be declared a value of $\psi_1$ or $\psi_2$.
To recover coordinates of those linear unknowns from
(B.43.17), one must subtract a specified reference origin
and its full quadratic pullback map.  The explicit covers
provide the actual affine target data.  That remaining
comparison, and the entire arbitrary-input image of B.5,
are not claimed to have been evaluated by the probes alone.
\end{proof}
